\chapter{Calibration}\label{chap:calibration}

Calibration computes the base value of every variable and every parameter of
the model from the SAM and from the elasticities and other settings given by
the user, so that the base solves every equation. \manref{cge_map} assigns
the accounts to roles, checks the SAM and rewrites the common
national-accounts layouts; \manref{cge_calibrate} then computes the
parameters and checks benchmark replication. In this chapter $\SAM(r, c)$ is
a cell of the SAM after the rewrites, and a superscript $0$ marks a base value.

\section{SAM checks and the flows the model represents}\label{sec:flows}

Before calibration the SAM must be square, with identical row and column
labels; it must be balanced, each row total equal to its column total within
$10^{-8}$ times the largest absolute cell; and every account with a non-zero
flow must have a role. Every non-zero cell must be one of the flows of the
table below. \manref{cge_map} lists every other flow, grouped by the roles of
the paying and the receiving accounts, before it stops.

\begin{cgetable}{p{0.24\linewidth}p{0.72\linewidth}}
\toprule
Paid by (column) & Received by (row)\\
\midrule\endhead
\bottomrule\endlastfoot
commodity & activities (domestic output); rest of the world (imports); import, sales, export and VAT accounts; margin accounts (margins)\\
activity & commodities (intermediate inputs); factors; activity tax account; accounts of taxes on factor use\\
factor & households; enterprises; government; rest of the world; health insurance fund (contributions); account of taxes on factor income\\
household & commodities; activities (home consumption); direct tax account; savings--investment; households and enterprises (transfers); rest of the world (transfers)\\
enterprise & direct tax account; savings--investment; households and enterprises (transfers); rest of the world (transfers)\\
government & commodities; households and enterprises (transfers); government (internal transfers); rest of the world (transfers); savings--investment; health insurance fund\\
tax account & government\\
savings--investment & commodities (investment, including changes in inventories)\\
margin account & commodities (margin services)\\
health insurance fund & commodities (purchases); savings--investment\\
rest of the world & commodities (exports); households, enterprises and government (transfers); factors (factor income from abroad); savings--investment; health insurance fund\\
\end{cgetable}

Negative cells are allowed only in tax rows, the savings--investment row,
transfer cells (the flows labelled transfers above, and factor income from
abroad), payments of tax accounts to the government, VAT cells, investment
demand (a fall in inventories) and sales of commodities by activities (a
negative entry of the make matrix gives a negative yield; each activity's
output and each commodity's domestic output must still be positive). Every
activity needs positive output and positive value added, and every commodity
needs domestic output or imports. An account with no flows at all needs no
role.

\section{Layouts rewritten before the checks}\label{sec:layouts}

National SAMs often record some flows differently from the layout of the
model. \manref{cge_map} rewrites the common ones before the checks above and
reports the amounts it moves; the SAM itself is not changed, and
\manref{cge_samexport} with \cmd{mapped} writes the SAM after the rewrites.

\subsection{Exports recorded as sales by activities}

A cell $\SAM(a, row)$ is treated as exports of the commodity paired with
$a$: it is added to $\SAM(c(a), row)$ and to $\SAM(a, c(a))$, and set to zero.
The activity's output and the rest of the world's spending are unchanged; the
commodity's supply and demand both rise by the exports.

\subsection{Taxes paid directly to the government}

A tax option of the mapping (direct, sales, import, export or activity tax)
may name the government account itself. The cells of the government's row in
the columns that pay that tax (households and enterprises for the direct tax,
commodities for the sales, import and export taxes, activities for the
activity tax) are then moved to an internal tax account, whose column pays
their total to the government. Government revenue is unchanged. At most one of
the sales, import and export tax options may name the government account,
because the cells would otherwise be ambiguous.

\subsection{Margins netted out of margin activities}

Some SAMs attribute trade and transport margins to the producing activities,
as their purchases of trade and transport services, and net the margin
revenue out of the output of the trade and transport activities by recording
a negative use of their own product, $\SAM(c(a), a) < 0$. For each activity
named in the option \cmd{margins()}, $-\SAM(c(a), a)$ is added to its output
$\SAM(a, c(a))$, and $\SAM(c(a), a)$ is set to zero. The activity's output
then equals its costs, and the margins enter the model as the producing
activities' fixed-coefficient purchases of trade and transport services, as
the SAM records them. Naming an activity without a negative use of its own
product is an error. Only the net amount is recorded, so after the rewrite the
activity's use of its own product is zero.

\subsection{Territorial correction}

National accounts move household consumption from the territorial to the
residents' concept with a correction account: residents' purchases abroad are
recorded as its imports, non-residents' purchases in the country as its
exports (sometimes in the row of a matching activity), and the net as a
negative household consumption of the correction. For the commodity (and,
optionally, the activity) named in \cmd{territorial()}, the activity's exports
are first moved to the commodity, and the activity must then have no flows.
With $IM$ the commodity's imports, $EX$ its exports and $HC_h$ household $h$'s
consumption of it, $\sum_h HC_h = IM - EX$, each household's share is
$s_h = HC_h / \sum_h HC_h$, and
\begin{equation*}
\SAM(row, h) \mathrel{+}= s_h\,IM, \qquad \SAM(h, row) \mathrel{+}= s_h\,EX:
\end{equation*}
residents' purchases abroad become household payments to the rest of the
world, and non-residents' purchases become transfers from the rest of the
world to households. The correction's cells are then set to zero.

\subsection{Stock-change accounts}

Some SAMs record changes in inventories in a separate account, financed by
savings--investment and buying commodities. For an account named in
\cmd{stocks()}, its payments are added to the savings--investment column, and
its row and column are set to zero. Changes in inventories are then part of
investment demand, scaled by $\V{IADJ}$ \eqref{eq:I6}.

\subsection{VAT within the sales tax}

When the SAM records VAT within the sales tax account, \cmd{vatshare()} gives
the VAT share of each commodity's sales tax; that share of
$\SAM(tax_{sal}, c)$ is moved to an internal VAT account before the checks
(\secref{sec:vat}). \cmd{vatshare(0)} creates an empty VAT account, so that a
VAT can be introduced by a shock.

\subsection{Roles that are not rewrites}

Margin accounts, mapped with \cmd{marginaccounts()} or, for margins by market,
with \cmd{domesticmargins()}, \cmd{importmargins()} and
\cmd{exportmargins()}, collect margins from commodities, $\SAM(t, c)$, and buy margin services,
$\SAM(c', t)$; they may have no other flows. A health insurance fund
(\cmd{insurance()}) receives contributions from factors and transfers from the
government and the rest of the world, buys health services and saves. A
household's payment to an activity, $\SAM(a, h)$, is home consumption of the
activity's output (\secref{sec:home}).

The rewrites of exports, of taxes paid to the government, of the territorial
correction and of stock changes leave every total that enters the model
unchanged, except that the territorial correction replaces the correction's
imports, exports and negative consumption by transfers. The margin rewrite
raises the output of the margin activities, and the supply and demand of
their products, by the margins.

\section{Normalisation}\label{sec:normalisation}

All base prices are one:
\begin{equation*}
\V{PA}^0_a = \V{PVA}^0_a = \V{PX}^0_c = \V{PD}^0_c = \V{PE}^0_c = \V{PQ}^0_c = 1,
\qquad \V{EXR}^0 = 1, \qquad \V{pwm}_c = 1,
\end{equation*}
and $\V{W}^0_f = \V{WFDIST}^0_{fa} = 1$ unless factor quantities are supplied
(\secref{sec:factor-quantities}). The base quantity of every flow is therefore
its value in the SAM.

\section{Taxes and trade}\label{sec:cal-trade}

\begin{itemize}
\item Imports: $\V{QM}^0_c = \SAM(row, c)$, $\V{tm}_c = \SAM(tax_{imp}, c) / \V{QM}^0_c$
and $\V{PM}^0_c = 1 + \V{tm}_c$.
\item Activities: $\V{QA}^0_a$ is the column total of activity $a$, and
$\V{ta}_a = \SAM(tax_{act}, a) / \V{QA}^0_a$.
\item Yields: $\theta_{ac} = \SAM(a, c) / \V{QA}^0_a$, and commodity output
$\V{QX}^0_c = \sum_a \SAM(a, c)$; $CX$ is the set of commodities with
$\V{QX}^0_c > 0$. With home consumption, marketed output replaces
$\V{QA}^0_a$ in the yields (\secref{sec:home}).
\item Exports: a commodity with $\SAM(c, row) > 0$ belongs to $CR$ if
$\SAM(c, row) - \SAM(tax_{exp}, c) \ge \V{QX}^0_c$, and to $CE$ otherwise.
\item For $c \in CE$: $\V{te}_c = \SAM(tax_{exp}, c) / \SAM(c, row)$,
$\V{QE}^0_c = \SAM(c, row) - \SAM(tax_{exp}, c)$, $\V{pwe}_c = 1/(1 - \V{te}_c)$
and $\V{QD}^0_c = \V{QX}^0_c - \V{QE}^0_c$.
\item For $c \in CR$: $\overline{\V{qer}}_c = \V{QER}^0_c = \SAM(c, row)$,
$\V{pwe}_c = 1$ and $\V{QD}^0_c = \V{QX}^0_c$. An export tax on such a commodity
is an error.
\item For other commodities, $\V{QD}^0_c = \V{QX}^0_c$.
\item Margins paid by commodity $c$: $\V{MG}^0_c = \sum_{t \in T} \SAM(t, c)$.
\item Composite supply and the sales tax rate:
\begin{equation*}
\V{QQ}^0_c = \V{QD}^0_c + \V{PM}^0_c\,\V{QM}^0_c + \SAM(tax_{sal}, c) + \V{MG}^0_c,
\qquad \V{tq}_c = \SAM(tax_{sal}, c) / \V{QQ}^0_c .
\end{equation*}
VAT is not part of
$\V{QQ}^0_c$ (\secref{sec:vat}).
\item Margin services: $\V{QT}^0_{c'} = \sum_{t \in T} \SAM(c', t)$ and
\begin{equation*}
\V{icm}_{c'c} = \frac{1}{\V{QQ}^0_c} \sum_{t \in T} \SAM(t, c)\,\frac{\SAM(c', t)}{\sum_{c''} \SAM(c'', t)},
\end{equation*}
so that the margins collected by account $t$ are spent on services in the
proportions in which $t$ buys them.
\end{itemize}
Balance of the SAM implies that $\V{QQ}^0_c$ equals total demand for $c$ at
purchaser prices other than exports from domestic output and VAT: the
commodity's row total less $\V{QE}^0_c$ and $\SAM(vat, c)$. Calibration checks
this identity rather than assuming it.

\section{Production}\label{sec:cal-production}

\begin{itemize}
\item Intermediate input: $\V{QINTA}^0_a = \sum_c \SAM(c, a)$ and
$\V{ica}_{ca} = \SAM(c, a) / \V{QINTA}^0_a$ (with VAT, the quantities of
\secref{sec:vat}).
\item Factors: $\V{QF}^0_{fa} = \SAM(f, a) / (\V{W}^0_f\,\V{WFDIST}^0_{fa})$ and
$\V{QVA}^0_a = \sum_f \SAM(f, a)$ (with taxes on factor use, their cost is
added; \secref{sec:factortax}).
\item Leontief top nest: $\V{iva}_a = \V{QVA}^0_a / \V{QA}^0_a$ and
$\V{inta}_a = \V{QINTA}^0_a / \V{QA}^0_a$.
\item CES top nest: with
$r_a = (\V{PVA}^0_a / \V{PINTA}^0_a)\,(\V{QVA}^0_a / \V{QINTA}^0_a)^{1 + \rho^a_a}$,
$\delta^a_a = r_a / (1 + r_a)$ and $1 - \delta^a_a = 1/(1 + r_a)$ (both
stored), and
\begin{equation*}
\V{aa}_a = \V{QA}^0_a \Big/ \Bigl[\delta^a_a\,(\V{QVA}^0_a)^{-\rho^a_a} + (1 - \delta^a_a)\,(\V{QINTA}^0_a)^{-\rho^a_a}\Bigr]^{-1/\rho^a_a}.
\end{equation*}
\item Cobb--Douglas top nest:
$\delta^a_a = \V{PVA}^0_a\,\V{QVA}^0_a / (\V{PVA}^0_a\,\V{QVA}^0_a + \V{PINTA}^0_a\,\V{QINTA}^0_a)$
and $\V{aa}_a = \V{QA}^0_a / \bigl[(\V{QVA}^0_a)^{\delta^a_a}(\V{QINTA}^0_a)^{1 - \delta^a_a}\bigr]$.
\item CES value added: with $w_{fa} = \V{W}^0_f\,\V{WFDIST}^0_{fa}$,
\begin{equation*}
\delta^{va}_{fa} = \frac{w_{fa}\,(\V{QF}^0_{fa})^{1 + \rho^{va}_a}}{\sum_{f'} w_{f'a}\,(\V{QF}^0_{f'a})^{1 + \rho^{va}_a}},
\qquad
\V{ad}_a = \V{QVA}^0_a \Big/ \Bigl[\sum_f \delta^{va}_{fa}\,(\V{QF}^0_{fa})^{-\rho^{va}_a}\Bigr]^{-1/\rho^{va}_a}.
\end{equation*}
\item Cobb--Douglas value added: $\delta^{va}_{fa} = \SAM(f, a) / \V{QVA}^0_a$ and
$\V{ad}_a = \V{QVA}^0_a / \prod_f (\V{QF}^0_{fa})^{\delta^{va}_{fa}}$.
\item Factor efficiency: $\V{af}_f = \V{EFF}^0_f = \V{DW}^0_f = 1$, so the
parameters above are unchanged, and
\begin{equation*}
\V{HS}^0 = \sum_{c \in HC} \Bigl[\sum_h \SAM(c, h) + \SAM(c, gov) + \SAM(c, shi)\Bigr],
\end{equation*}
which must be positive when some health elasticity is positive.
\end{itemize}

\section{CET and Armington functions}\label{sec:cal-cet}

\begin{itemize}
\item CET:
\begin{gather*}
\delta^t_c = \frac{1}{1 + (\V{PD}^0_c / \V{PE}^0_c)\,(\V{QE}^0_c / \V{QD}^0_c)^{\rho^t_c - 1}},\\
\V{at}_c = \V{QX}^0_c \Big/ \Bigl[\delta^t_c\,(\V{QE}^0_c)^{\rho^t_c} + (1 - \delta^t_c)\,(\V{QD}^0_c)^{\rho^t_c}\Bigr]^{1/\rho^t_c}.
\end{gather*}
\item Armington, for imported commodities with domestic output:
\begin{gather*}
\delta^q_c = \frac{(\V{PM}^0_c / \V{PD}^0_c)\,(\V{QM}^0_c / \V{QD}^0_c)^{1 + \rho^q_c}}{1 + (\V{PM}^0_c / \V{PD}^0_c)\,(\V{QM}^0_c / \V{QD}^0_c)^{1 + \rho^q_c}},\\
\V{aq}_c = \V{QQ}^0_c \Big/ \Bigl[\delta^q_c\,(\V{QM}^0_c)^{-\rho^q_c} + (1 - \delta^q_c)\,(\V{QD}^0_c)^{-\rho^q_c}\Bigr]^{-1/\rho^q_c};
\end{gather*}
$\V{aq}_c = \V{QQ}^0_c / \V{QD}^0_c$ for commodities that are not imported, and
$\V{aq}_c = \V{QQ}^0_c / \V{QM}^0_c$ for those that are imported only. With
margins by market, $\V{PDD}^0_c$ and $\V{PM}^0_c$ include the margins
(\secref{sec:margins-market}).
\end{itemize}
The complements $1 - \delta^t_c$ and $1 - \delta^q_c$ are stored as separate
parameters, computed from the same ratio as the shares (for example
$1 - \delta^t_c = r/(1 + r)$ with $\delta^t_c = 1/(1 + r)$), and used in
\eqref{eq:T1}, \eqref{eq:T2}, \eqref{eq:T4} and \eqref{eq:T5}. Computing
$1 - \delta$ by subtraction loses most significant digits when a share is
close to one, as with a low CET elasticity and a small export share, and can
leave a benchmark residual near $10^{-7}$.

\section{Institutions}\label{sec:cal-institutions}

\begin{itemize}
\item Factor income: $\V{YF}^0_f$ is the row total of factor $f$, and
\begin{equation*}
\V{tshi}_f = \frac{\SAM(shi, f)}{\V{YF}^0_f}, \qquad
\V{shif}_{if} = \frac{\SAM(i, f)}{(1 - \V{tshi}_f)\,\V{YF}^0_f}, \quad i \ne shi,
\end{equation*}
with $\V{tshi}_f = 0$ without a fund; with a tax on factor income the
denominator is $(1 - \V{tshi}_f - \V{tfy}_f)\,\V{YF}^0_f$.
\item Health insurance fund: $\V{YSHI}^0$ is the row total of $shi$, the
base purchases are $\overline{\V{qshi}}_c = \SAM(c, shi)$ and the savings rate
is $\V{sshi} = \SAM(sav, shi) / \V{YSHI}^0$; the transfers are
$\V{trnsfr}_{shi,gov} = \SAM(shi, gov)$ and
$\V{trnsfr}_{shi,row} = \SAM(shi, row)$; and $\V{SHIADJ}^0 = 1$.
\item Domestic institutions: $\V{YI}^0_d$ is the row total of $d \in D$;
$\overline{\V{tins}}_d = \SAM(tax_{dir}, d) / \V{YI}^0_d$ and
$\overline{\V{mps}}_d = \SAM(sav, d) / \bigl[(1 - \overline{\V{tins}}_d)\,\V{YI}^0_d\bigr]$.
\item Income after direct taxes and savings:
$\V{DY}^0_d = \V{YI}^0_d - \SAM(tax_{dir}, d) - \SAM(sav, d)$; calibration stops
with an error if $\V{DY}^0_d = 0$ while $d$ makes transfers or consumes. Then
$\V{shii}_{dd'} = \SAM(d, d') / \V{DY}^0_{d'}$ and
$\V{shir}_d = \SAM(row, d) / \V{DY}^0_d$.
\item Consumption: $\V{EH}^0_h = \sum_c \SAM(c, h)$ (plus home consumption),
$\V{hc}_h = \V{EH}^0_h / \V{DY}^0_h$, and budget shares
$w_{ch} = \SAM(c, h) / \V{EH}^0_h$.
\item Cobb--Douglas demand, the default: $\beta_{ch} = w_{ch}$ and
$\gamma_{ch} = 0$.
\item Linear expenditure system: given income elasticities
$\varepsilon_c > 0$ and a Frisch parameter $\phi_h < -1$ for each household,
\begin{equation*}
\beta_{ch} = \frac{\varepsilon_c\,w_{ch}}{\sum_{c'} \varepsilon_{c'}\,w_{c'h}},
\qquad
\gamma_{ch} = \V{QH}^0_{ch} + \frac{\beta_{ch}\,\V{EH}^0_h}{\V{PH}^0_c\,\phi_h},
\end{equation*}
so that the marginal shares sum to one (the income elasticities used are then
$\beta_{ch}/w_{ch}$, proportional to those supplied). Supernumerary spending,
$\V{EH}^0_h - \sum_c \V{PH}^0_c\,\gamma_{ch}$, is then $-\V{EH}^0_h/\phi_h$,
positive and below total spending because $\phi_h < -1$.
\item Transfers: $\V{trnsfr}_{d,gov} = \SAM(d, gov)$,
$\V{trnsfr}_{gov,gov} = \SAM(gov, gov)$, $\V{trnsfr}_{d,row} = \SAM(d, row)$,
$\V{trnsfr}_{f,row} = \SAM(f, row)$, $\V{trnsfr}_{gov,row} = \SAM(gov, row)$ and
$\V{trnsfr}_{row,gov} = \SAM(row, gov)$ (with public debt, the interest is
separated from them; \secref{sec:debt}).
\item Final demand: $\overline{\V{qg}}_c = \SAM(c, gov)$ and
$\overline{\V{qinv}}_c = \SAM(c, sav)$ (divided by $1 + \V{tv}_c$ with VAT),
$\V{GSAV}^0 = \SAM(sav, gov)$ and $\V{FSAV}^0 = \SAM(sav, row)$.
\item Scaling factors: $\V{TINSADJ}^0 = \V{MPSADJ}^0 = \V{IADJ}^0 = \V{GADJ}^0 = 1$,
and the tax scaling factors are one.
\item Price index weights:
\begin{equation*}
\V{cwts}_c = \frac{\sum_h \SAM(c, h)}{(1 + \V{tv}_c) \sum_{c'} \sum_h \SAM(c', h)},
\qquad
\V{dwts}_c = \frac{\V{QD}^0_c}{\sum_{c'} \V{QD}^0_{c'}},
\end{equation*}
so that $\V{CPI}^0 = \V{DPI}^0 = 1$.
\end{itemize}
Balance of the SAM implies $\V{hc}_h = 1 - \sum_d \V{shii}_{dh} - \V{shir}_h$
for households and $\sum_d \V{shii}_{de} + \V{shir}_e = 1$ for enterprises;
calibration checks both.

\section{Tax accounts}

Each tax is recorded in a tax account whose row collects the tax and whose
column pays only the government, either in the SAM itself or as an internal
account created by a rewrite (\secref{sec:layouts}).

\section{Factor quantities}\label{sec:factor-quantities}

If the user supplies base factor quantities $\V{QF}^0_{fa}$, for example
employment, the base prices of the factors are
\begin{equation*}
\V{W}^0_f = \frac{\sum_a \SAM(f, a)}{\sum_a \V{QF}^0_{fa}},
\qquad
\V{WFDIST}^0_{fa} = \frac{\SAM(f, a)}{\V{W}^0_f\,\V{QF}^0_{fa}},
\end{equation*}
so that the average price of a factor is its payment per unit, and the wage
distortions are the differences between activities. Otherwise
$\V{W}^0_f = \V{WFDIST}^0_{fa} = 1$.

\section{Benchmark replication}

The calibrated base must solve every equation of the model. \manref{cge_calibrate}
evaluates every scaled residual at the base, displays the largest and stops
with an error if it exceeds $10^{-9}$. A failure points to a problem in the
SAM or the mapping, such as a negative cell where values must be
non-negative or an elasticity that the model does not allow, not to the
solver. With the default closure, the solver started from the base returns it
in zero iterations; chapter~\ref{chap:checks} lists the further checks.
